\documentclass[12pt]{article}

\usepackage[a4paper,top=2.0cm,left=2.5cm,right=2.5cm,bottom=2cm,includefoot,includehead]{geometry}
\usepackage[serbian]{babel}

\input glyphtounicode
\pdfgentounicode=1

\usepackage[noslant]{srbtiks}
\usepackage[intlimits]{amsmath}
\usepackage{fonttable}
\usepackage{fancyvrb}

\title{Paket \lat{\textbf{srbtiks}}}
\author{\textsc{Uroš Stefanović}\footnote{\lat{urostajms@gmail.com}}}

\makeatletter
\renewcommand{\@seccntformat}[1]{\csname the#1\endcsname.\quad}
\makeatother

\begin{document}

\maketitle

\section{Šta je \lat{\emph{srbtiks}}?}

Paket \textbf{\lat{srbtiks}} omogućava korišćenje fonta \textbf{\lat{stix2}} za srpski i makedonski jezik za \lat{\TeX{} (\LaTeX)}.
Prema tome, da bi paket pravilno radio, potrebno je da paket \textbf{\lat{stix2-type1} }bude instaliran.
Novina je to što je za ćirilička slova upotrebljeno kodiranje \emph{\lat{T1}} umesto \emph{\lat{T2A}} kodiranja;
time je postignut efekat da se ćirilička slova unose preko latiničke tastature.

Namena ovog paketa je da omogući jednostavno preslovljavanje iz latinice u ćirilicu.
Primetno je da veliki broj korisnika \lat{\LaTeX}"-a preferira latinički unos
(što je često jednostavnije zbog velikog broja matematčikih formula),
pa se često pribegava korišćenju \emph{\lat{OT2}} kodiranja, što je veoma loše rešenje.
Cilj paketa \textbf{\lat{srbtiks}} je da pruži kvalitetnu alternativu svim postojećim paketima koji nude mogućnost preslovljavanja u ćirilicu, takođe je fenomenalna alternativa fontu \emph{\lat{Times New Roman}}.

\section{Kako se koristi?}

Paket se koristi jednostavnim pozivom \textbf{\lat{\textbackslash usepackage\{srbtiks\}}}.
Prilikom pozivanja automatski će se učitati paket \textbf{\lat{stix2}} sa opcijom \textbf{\lat{upint}}.
Ako ne želimo da se taj paket učita onda možemo da iskoristimo opciju \textbf{\lat{\textbackslash usepackage[nostix]\{srbtiks\}}}.

Font \textbf{\lat{srbtiks}} podržava samo \emph{\lat{T1}} i \emph{\lat{TS1}} kodiranja, tako da za latinički test potrebno je koristiti font \textbf{\lat{stix2}}.
Zbog toga su dostupne komande \textbf{\lat{\textbackslash LAT}}, \textbf{\lat{\textbackslash CYR}}, \textbf{\lat{\textbackslash lat}} i \textbf{\lat{\textbackslash cyr}} koje omogućavaju prelaz sa jednog fonta na drugi.
Komande \textbf{\lat{\textbackslash LAT}} i \textbf{\lat{\textbackslash CYR}} se koriste bez argumenta dok \textbf{\lat{\textbackslash lat}} i \textbf{\lat{\textbackslash cyr}} imaju jedan argument, na primer \textbf{\lat{\textbackslash lat\{ovaj tekst je ispisan latinicom\}}}.

\section{Unos}

Slova je moguće unositi na nekoliko načina: može da se koristi \textit{\lat{utf8}} kodiranje, pa tako \textbf{\lat{š đ č ć ž dž}} daju redom \textbf{š đ č ć ž dž}.
Drugi način je preko kombinacije akcenata, na primer \textbf{\lat{\textbackslash v~s \textbackslash dj \textbackslash v~z}} daju \textbf{\v s \dj{} \v z}.
Treći način je kombinacijom sa slovima \textbf{\lat{x}} i \textbf{\lat{y}}, tako imamo da \textbf{\lat{sx dx dy zx cx cy ny ly}} daju \textbf{sx dx dy zx cx cy ny ly} (\emph{detaljnije u tabeli~\ref{t1} na strani~\pageref{t1}}).

\section{Specijalni znaci}

Ćirilička slova koja se ne koriste u srpskom jeziku (na primer \char215{} ili \char197) su i dalje dostupna.
Na strani~\pageref{t2} mogu da se vide svi dostupni karakteri. Jedan od načina da se pristupi znacima je preko komande \textbf{\lat{\textbackslash char\#}}, gde je \textbf{\#} redni broj karaktera; tako na primer \textbf{\lat{\textbackslash char179}} daje \char179.

Akcenti koji se koriste u srpskom jezku su \emph{kratkosilazni}, \emph{kratkouzlazni}, \emph{dugosilazni} i \emph{dugouzlazni}.
Takođe su neophodni \emph{znak dužine} i \emph{znak kratkoće}, kao i \emph{genitivni znak}.
Svi su dostupni u ovom paketu, pa tako komande (u kombinaciji sa slovom \textbf{e}) \textbf{\lat{\textbackslash AKS}}, \textbf{\lat{\textbackslash AKU}}, \textbf{\lat{\textbackslash ADS}}, \textbf{\lat{\textbackslash ADU}}, \textbf{\lat{\textbackslash ADZ}}, \textbf{\lat{\textbackslash AKZ}} i \textbf{\lat{\textbackslash AGZ}} daju redom \textbf{\AKS e}, \textbf{\AKU e}, \textbf{\ADS e}, \textbf{\ADU e}, \textbf{\ADZ e}, \textbf{\AKZ e} i \textbf{\AGZ e}.

Znaci poput \textnumero, §, ÷, × su dostupni preko kodiranja \textit{\lat{TS1}}.
Za više detalja pogledati dokumentaciju paketa \textbf{\lat{stix2-type1}}.

\section{Matematika}

Kada se paket učita komande \textbf{\lat{\textbackslash leq}}, \textbf{\lat{\textbackslash geq}}, \textbf{\lat{\textbackslash nleq}} i \textbf{\lat{\textbackslash ngeq}} će redom dati znake $\leqslant$, $\geqslant$, $\nleqslant$ i $\ngeqslant$ umesto standardnih $\leq$, $\geq$, $\nleq$ i $\ngeq$ (ali samo ako nije upotrebljena opcija \textbf{\lat{nostix}}!).
Ovo može da se spreči korišćenjem opcije \textbf{\lat{noslant}} prilikom učitavanja paketa.


\section{Problemi}

Kako se ne bi desilo da se rimski brojevi ispišu kao na primer II ili XVII,
komande \lat{\textbf{\textbackslash@roman}} i \lat{\textbf{\textbackslash@Roman}} su izmenjene.
\emph{Zbog toga postoji mogućnost da se rimski brojevi ne prikažu uvek kako je zamišljeno.}

Takođe treba obratiti pažnju da se neke komande, na primer \lat{\textbf{\textbackslash LaTeX}}, uvek pišu kao \lat{\textbf{\textbackslash lat\{\textbackslash LaTeX\}}} kako se ne bi prikazale kao \LaTeX.

Sav tekst se unosi isključivo latinicom, ćirilički unos trenutno nije moguć.
Zato je bolji izbor korišćenje opcije \textbf{\lat{serbian}} umesto \textbf{\lat{serbianc}} za paket \textbf{\lat{babel}}.

Akcente \kern.5em\char28\kern.6em i \kern.5em\char29\kern.6em, ako pišemo kao izolovane znake, treba pisati recimo komandama \textbf{\lat{\textbackslash kern.5em\textbackslash char28\textbackslash kern.6em}} i \textbf{\lat{\textbackslash kern.5em\textbackslash char29\textbackslash kern.6em}} kako bi se pravilno prikazali.

\section{Napomena}
\emph{Ovaj paket treba posmatrati kao eksperimenalni rad.
Promene na paketu su moguće ako se paket \textbf{\lat{stix2-type1}} promeni.
Takođe postoji mogućnost da ovaj paket ne prati eventualnu promenu paketa \textbf{\lat{stix2-type1}}.}

\pagebreak

\section{Primer}

\textbf{Ulaz:}
\begin{Verbatim}
\documentclass[12pt]{article}
\usepackage{srbtiks}
\usepackage[serbian]{babel}
\usepackage[intlimits]{amsmath}
\begin{document}
Ovo je tekst napisan ćirilicom.
Mozxe da se pisxe ovako, a mo\v ze i ovako.
Kracyi tekst latinicom se pisxe \lat{ovako}.
\LAT Ovako se piše duži tekst latinicom,
\CYR a ovako se vraća na ćirilicu.
\textit{abvgddygyezxzzyijkllymnnyoprstcykyufhccxdxsx}

Matematičke formule: $x^2+y^2=z^2$,
ili $\sum_{i=1}^n=\frac{n(n+1)}{2}$.
Može i
\[ \int_{-\infty}^\infty x\,\mathrm{d}x + \prod_0^4f(\vec{z}) \leq\pi. \]
\end{document}
\end{Verbatim}

\textbf{Izlaz:}

Ovo je tekst napisan ćirilicom.
Mozxe da se pisxe ovako, a mo\v ze i ovako.
Kracyi tekst latinicom se pisxe \lat{ovako}.
\LAT Ovako se piše duži tekst latinicom,
\CYR a ovako se vraća na ćirilicu.
\textit{abvgddygyezxzzyijkllymnnyoprstcykyufhccxdxsx}

Matematičke formule: $x^2+y^2=z^2$,
ili $\sum_{i=1}^n=\frac{n(n+1)}{2}$.
Može i
\[ \int_{-\infty}^\infty x\,\mathrm{d}x + \prod_0^4f(\vec{z}) \leqslant\pi. \]

\pagebreak

\begin{table}[h]
\centering
\begin{tabular}{|cc|cc|cc|cc|cc|}
\hline
Ulaz & Izlaz & Ulaz & Izlaz & Ulaz & Izlaz & Ulaz & Izlaz & Ulaz & Izlaz  \\
\hline\hline
\lat{a} & a & \lat{b} & b & \lat{c} & c & \lat{d} & d & \lat{e} & e \\
\hline
\lat{f} & f & \lat{g} & g & \lat{h} & h & \lat{i} & i & \lat{j} & j \\
\hline
\lat{k} & k & \lat{l} & l & \lat{m} & m & \lat{n} & n & \lat{o} & o \\
\hline
\lat{p} & p & \lat{q} & q & \lat{r} & r & \lat{s} & s & \lat{t} & t \\
\hline
\lat{u} & u & \lat{v} & v & \lat{w} & w & \lat{x} & x & \lat{y} & y \\
\hline
\lat{z} & z & \lat{cx} & cx & \lat{cy} & cy & \lat{dx} & dx & \lat{dy} & dy \\
\hline
\lat{ex} & ex & \lat{gy} & gy & \lat{ky} & ky & \lat{lj} & lj & \lat{ly} & ly \\
\hline
\lat{nj} & nj & \lat{ny} & ny & \lat{sx} & sx & \lat{zx} & zx & \lat{zy} & zy \\
\hline
\lat{č} & č & \lat{ć} & ć & \lat{đ} & đ & \lat{š} & š & \lat{ž} & ž \\
\hline
\lat{dž} & dž & \lat{\textbackslash v c} & \v c & \lat{\textbackslash'c} & \'c & \lat{\textbackslash dj} & \dj & \lat{\textbackslash v e} & \v e \\
\hline
\lat{\textbackslash v s} & \v s & \lat{\textbackslash v z} & \v z & \lat{\textbackslash v l} & \v l & \lat{\textbackslash v n} & \v n & \lat{\textbackslash v d} & \v d \\
\hline
\lat{d\textbackslash v z} & d\v z & \lat{\textbackslash`e} & \`e & \lat{\textbackslash`i} & \`i & \lat{dzx} & dzx & \lat{Dzx} & Dzx \\
\hline\hline
\lat{A} & A & \lat{B} & B & \lat{C} & C & \lat{D} & D & \lat{E} & E \\
\hline
\lat{F} & F & \lat{G} & G & \lat{H} & H & \lat{I} & I & \lat{J} & J \\
\hline
\lat{K} & K & \lat{L} & L & \lat{M} & M & \lat{N} & N & \lat{O} & O \\
\hline
\lat{P} & P & \lat{Q} & Q & \lat{R} & R & \lat{S} & S & \lat{T} & T \\
\hline
\lat{U} & U & \lat{V} & V & \lat{W} & W & \lat{X} & X & \lat{Y} & Y \\
\hline
\lat{Z} & Z & \lat{CX} & CX & \lat{CY} & CY & \lat{DX} & DX & \lat{DY} & DY \\
\hline
\lat{EX} & EX & \lat{GY} & GY & \lat{KY} & KY & \lat{LJ} & LJ & \lat{LY} & LY \\
\hline
\lat{NJ} & NJ & \lat{NY} & NY & \lat{SX} & SX & \lat{ZX} & ZX & \lat{ZY} & ZY \\
\hline
\lat{Č} & Č & \lat{Ć} & Ć & \lat{Đ} & Đ & \lat{Š} & Š & \lat{Ž} & Ž \\
\hline
\lat{DŽ} & DŽ & \lat{\textbackslash v C} & \v C & \lat{\textbackslash'C} & \'C & \lat{\textbackslash DJ} & \DJ & \lat{\textbackslash v E} & \v E \\
\hline
\lat{\textbackslash v S} & \v S & \lat{\textbackslash v Z} & \v Z & \lat{\textbackslash v L} & \v L & \lat{\textbackslash v N} & \v N & \lat{\textbackslash v D} & \v D \\
\hline
\lat{D\textbackslash v Z} & D\v Z & \lat{\textbackslash`E} & \`E & \lat{\textbackslash`I} & \`I & \lat{D\textbackslash v z} & D\v z & \lat{DZX} & DZX \\
\hline\hline
\lat{Dž} & Dž  & \lat{Cx} & Cx & \lat{Cy} & Cy & \lat{Dx} & Dx & \lat{Dy} & Dy \\
\hline
\lat{Ex} & Ex & \lat{Gy} & Gy & \lat{Ky} & Ky & \lat{Lj} & Lj & \lat{Ly} & Ly \\
\hline
\lat{Nj} & Nj & \lat{Ny} & Ny & \lat{Sx} & Sx & \lat{Zx} & Zx & \lat{Zy} & Zy \\
\hline
\end{tabular}
\caption{\emph{Ulaz i odgovarajući izlaz}}\label{t1}
\end{table}

\pagebreak

\fonttable{t1srbtiksr}\label{t2}

\end{document}